\chapter{Diskussion}
Dieses Kapitel bewertet die Umsetzung der Anforderungen und ordnet die Ergebnisse kritisch ein. Dabei werden sowohl die erreichten Ziele als auch die offenen Punkte und Grenzen des Prototyps zusammengefasst.

\section{Erfüllung der Anforderungen}
Die wesentlichen Anforderungen an die Hardware wurden erreicht. Der Aufbau arbeitet im aktuellen Entwicklungsstand ordnungsgemäß, die Spannungsversorgung ist stabil und die grundlegenden Funktionen sind nachweisbar umgesetzt. Gleichzeitig zeigen sich noch kleinere Kommunikationsschwierigkeiten mit verschiedenen Tachographen, die sich in instabilen Sitzungen und vereinzelten Abbrüchen äußern und weiter adressiert werden müssen.

\section{Grenzen des Prototyps}
Der Prototyp ist auf einen technischen Nachweis ausgelegt und noch nicht für den flächendeckenden Feldeinsatz qualifiziert. Neben den genannten Kommunikationsschwierigkeiten mit einzelnen Tachographen fehlen umfangreiche Langzeittests unter realen Fahrzeugbedingungen sowie eine vollständige Abdeckung unterschiedlicher Hersteller- und Gerätegenerationen. Damit ist die Übertragbarkeit auf alle Einsatzszenarien derzeit eingeschränkt und erfordert weitere Validierung.
